\section{Introduction}
\label{sec:intro}

Security claims about machine learning appear faster than the field can compare
them. Arp et al.\ read \DDpapers{} papers from CCS, S\&P, USENIX Security and NDSS
across a decade, name \DDpitfalls{} recurring pitfalls, and find every one of the
\DDpapers{} papers affected by at least \DDminPitfalls{} of them \cite{AQPWPWCR22}.
One of the \DDpitfalls{} is the threat model. Their measurement is of practice, and
the same defect shows up in the results themselves wherever one attack has been
studied by several groups.

Membership inference makes the defect concrete, and we use it to state the three
problems in Section~\ref{sec:problem}. The attack decides whether one record was in
a model's training set, and Meeus et al.\ count at least \MIAnewMethods{} new
methods for it proposed in 2024 alone \cite{MSJFRM25}. The published verdicts on whether it works span the full range.
Carlini et al.\ report that it succeeds at low false positive rates
\cite{CCNSTT22}. Duan et al.\ run five attacks against models whose training set is
known and report that none exceeds \MIAmimirAUC{} AUC outside one domain
\cite{DSMMSZTCEH24}. Das et al.\ report that on \MIAblindSets{} published
benchmarks a classifier that never queries the model beats the best attack that
does \cite{DZT25}. Zhang et al.\ argue that the attack cannot establish training-set
membership for a deployed model at all \cite{ZDKT25}. Every one of these results is
correct on its own terms. They differ in what they let the attacker know and spend,
and none of them states that in a form the others can be checked against.

Membership inference is one attack among many, and it appears in \NumMIApages{} of
the \NumSources{} works surveyed here. The corpus also covers evasion, poisoning,
backdoors, data reconstruction and model extraction, and on the other side
computational differential privacy, statistical-query and low-degree hardness,
bounded-memory adversaries and fine-grained cryptography. No single attack runs
through every part of the survey, so each part of Section~\ref{sec:landscape}
carries its own worked case.

The field has not ignored the adversary. It has described the adversary without
pricing it. The standing systematization of the area sorts attacks by the stage of
the pipeline the adversary acts at, whether it can write as well as read, the
access it holds, and the goal it pursues \cite{PMSW18}, and none of those four axes
is a computational resource. Three later systematizations, all of them preprints at the time
of writing, report that their own subfield lacks a formal definition of attack
success, on adversarial examples \cite{FK24}, on model extraction \cite{SSEP24},
and on data reconstruction \cite{WLBZ25}. Each proposes a repair, and each repair supplies knowledge, access
or a success metric. None supplies a budget. A reader holding two results therefore
cannot ask which adversary was stronger, because neither result answers it.

Cryptography has priced adversaries for fifty years, and the objects the comparison
needs already exist there. An adversary class fixes what the attacker may run. A
budget fixes how much of it the attacker may spend. A success criterion with a
tolerance fixes what counts as winning. A baseline fixes what public knowledge
alone already achieves, so that an advantage is measured against something. We
carry those four objects into machine learning security and use them to read the
field.

We read \NumSources{} formal works and record, for each one, which party's
resources enter its main claim and in what unit. Recording the budget per party
rather than per adversary is what lets a single statement hold a bound on the
attacker, a bound on the learner, a bound on an auditor, and a ratio between two of
them. It also makes the most common answer visible, which is that no party's
resources enter the claim at all.

Our contributions are the following.

\begin{itemize}
  \item A survey of \NumSources{} formal works on the security and privacy of
        learning, organized by the attack family together with the party each work
        constrains, recording for every work which party's resources enter its main
        claim, in what unit, and whether the work assesses the adversary explicitly
        or leaves the assessment implicit (Section~\ref{sec:landscape}).
  \item An experiment, $\mathrm{Exp}[\Pi, L, A, R, \kappa]$, whose four coordinates
        are the adversary's observation point on the chain from training data to
        model output, the criterion $R$ it must satisfy and the tolerance $\eta$ it
        is held to, a capability vector $\kappa$ carried per party, and the
        baseline. We show which published results each coordinate separates, and
        which results the experiment does not hold
        (Section~\ref{sec:formalization}).
  \item A mapping of the corpus into that experiment, generated from the recorded
        claim quantifier, antecedent and adversary bound of each work rather than
        assembled by hand, so that it can be rebuilt when the corpus changes
        (Section~\ref{sec:mapping}).
  \item Four regions of the experiment that no work in the corpus occupies, each
        reported with the search that established it rather than as an assertion,
        together with the count of works that state no budget for any party
        (Section~\ref{sec:findings}).
\end{itemize}
